The functions associated with the poly-mod section are very powerful. They provide tools for non-linear modulation: cross-modulation of the two voice oscillators and modulation of the filter cut-off frequency by VCO B. 

A familiar and common synthesizer feature is that oscillators and/or filters (frequency, filter cut-off frequency, pulse width, etc.) are modulated by a \textit{low frequency oscillator (LFO)}. An LFO does not operate in the audible frequency range (it is "low frequency") and it has a fixed frequency which is independent from the note played. Poly-mod is different in that the modulation itself is in the audible range and the resulting non-linearity changes the frequency spectrum drastically. The resulting effect is comparable to FM synthesis known from digital FM synthesizers\footnote{FM synthesizers typically use phase modulation. However, the term \textit{FM synthesis} is commonly used for the family of synthesis methods or modulations in which an audio range signal is used to modulate the frequency control.}, although in the Prophet-600 it is more raw and fully analog. Like the oscillator sync function, cross-modulation is hardware implemented in the Prophet-600. To be precise, the Prophet-600 implements \textit{exponential} FM by feeding the audio signal directly into the voltage control of the oscillator.

The poly-mod control can be found in the poly-mod sub-panel. 

\begin{center}
\scalebox{0.4}{
  \begin{tikzpicture}[scale=0.8]

    \addpar{-18cm,3.5cm}{\envrouting \\ The routing options determine whether poly-mod is driven by the filter or the 2nd envelope};
    \draw[line width = 2pt, dashed] (-4cm,3.5cm) -- (1cm,3.5cm);


   \polymodesubpanel{0,0}
  \end{tikzpicture}
}
\end{center}

There are two modulation sources, \polymodenv and \polymodosc and two destinations, the pitch of oscillator A  and the filter cut-off frequency. Each source is activated using the respective panel switch, \polymodfreq and \polymodfilter. The \polymodosc determines how strongly the oscillator B modulates the poly-mod targets (destinations). It is instructive to consider the effect of oscillator B modulation on the two targets first and then explore the effect of envelope modulation.

\textbf{Cross-modulation}

\begin{itemize}
  \item When you set the \polymodfreq switch to \textit{on}, oscillator B modulates the frequency of oscillator A (frequency modulation, FM). Depending on the modulation strength the result is a wide and rich frequency spectrum of oscillator A which contains among others the sum and difference of multiples of the two oscillators frequencies. The spectrum strongly depends on the frequency difference and may be totally non-harmonic but also harmonic for carefully adjusted frequency settings. Typical spectra are often described as bell sounding. Note that in contrast to digital FM synthesizers which use highly accurate sin waves, the oscillator output of the Prophet-600 does not have the same basic accuracy and the available waveforms already contain higher harmonics (triangle, sawtooth, etc.). Therefore, the FM result in the Prophet-600 is typically more drastic and dirtier and requires classic subtractive methods for shaping the sound, i.e. the application of the low pass filter. Another important difference to FM synthesizers is that the Prophet-600 (like most analog poly synthesizers with cross-modulation capability) uses exponential FM. It is characteristic of this feature that the modulation strength changes the pitch of the sound. If you want a harmonic patch for chord playing, you should treat \polymodosc as a static parameter once it is adjusted.
  \item When you set the \polymodfilter switch to \textit{on}, oscillator B modulates the filter cut-off frequency of the low pass filter. Like oscillator FM this is a non-linear modulation. This effect can also produce a wide frequency spectrum which may or may not be harmonic, depending on the relative frequency of oscillators A and B. The effect is more pronounced at low filter cut-off frequencies. The resonance also has a strong influence on the effect.  
\end{itemize}

Note that unless you use oscillator sync at the same time, poly-mod is very pitch sensitive so that patches which use strong poly-mod may need to be fine tuned using \freqfine or \oscfreq (for this, the \textit{free} option of the parameter \oscpitchmode is useful). Also, poly-mod patches may sound off on a cold unit due to temperature dependent drift and generally it maybe necessary to retune the synth more often when using extreme poly-mod patches. 

You should also be aware that the harmonic integrity of a patch may vary from instrument to instrument and - more importantly - that there can be voice-to-voice differences. Therefore, medium to strong poly-mod patches that sound good on one instrument may not immediately sound good on another unit or may even not be fully reproducible. In most cases you can compensate by adjusting \oscfreq (better A than B), \freqfine and/or \polymodosc. However, you may find that in your unit the differences among the voices are so strong that beyond a certain poly-mod strength, patches start sounding detuned (for example in chord play). Please refer to section \ref{fm} for more background on the design implications of poly-mod in the Prophet-600 and for measures you can take to reduce any detuning effects of poly-mod. Measures include scaling adjustment and the avoidance of the oscillator B pulse waveform in medium to strong poly-mod patches. 

Note that when you use oscillator sync, pitch related problems of poly-mod are not an issue.

\textbf{Envelope pitch modulation}

The second poly-mod modulation source is the poly-mod envelope. When the \polymodfreq switch is \textit{on}, the envelope modulates the pitch of oscillator A at the strength set by the \polymodenv dial. The \polymodenv control is bipolar. It  provides negative (inverse modulation, lowering the pitch) and positive modulation (raising the pitch). The pitch shift is harmonic and in the sustain phase of the envelope the maximum obtainable pitch shift (i.e. maximum positive or negative value of \polymodenv and sustain at full value) corresponds to $\pm 30$ semitones.

Note that the poly-mod envelope only affects the frequency of oscillator A (and never the filter cut-off frequency) even if the second destination, the \polymodfilter switch, is also \textit{on}. The filter cut-off frequency is already modulated by the filter envelope. Depending on the menu parameter \envrouting the source envelope is either the filter envelope or the 2nd envelope, see below for details on the options.

Pitch modulation by envelope is not actually a cross-modulation function. However, the association with poly-mod is useful because of the strong effect pitch (difference) has on cross-modulation. What can it be used for? Envelope modulation is a form of time modulation. It can change the sound from attack through hold to release. By time modulating the pitch in combination with cross-modulation (and/or sync), the sound can be made to sweep through different spectra, ranging from apocalyptic noise through bell sounds to delicate transients. The way humans perceive natural sounds is strongly influenced by the frequency characteristics of the transients (typically short lived high frequency components) in the attack and decay phases of the sound. This is best modelled using an envelope. So, being able to modulate the composition of frequency spectrum by envelope is a versatile tool for advanced sound design. This was probably the original thought behind putting the \polymodenv control in the poly-mod section.

In order to further expand the sound design possibilities, the latest firmware upgrade offers an envelope routing option, which is inspired by a similar setup in the Roland JX8P synthesizer. Having different envelopes modulating the filter and the poly-mod opens up new possibilities in the sound and spectrum dynamics. The following section explains the routing options. 

The menu parameter \envrouting determines if the filter envelope or the 2nd envelope is used in poly-mod\footnote{Up to version 2.1 RC3 the filter envelope was "hard wired" to poly-mod.}.

\begin{itemize}
  \setlength\itemsep{0cm}
  \item \textit{standard}: In this "standard" routing the amplitude is modulated by the 2nd envelope and the poly-mod is modulated by the filter envelope. This is the setup of the original Prophet-600 and versions 2.0 and 2.1 RC3. It is also the default in \version.
  \item \textit{poly-amplitude}: In this routing both the amplitude and the poly-mod are modulated by the 2nd envelope. 
  \item \textit{poly}: In this routing the amplitude is modulated by the filter envelope and the poly-mod is modulated exclusively by the 2nd envelope. Therefore the poly-mod time modulation is completely free, at the expense of having filter modulation and amplitude tied together.
  \item \textit{gate}: In this routing the amplitude is not modulated at all, but behaves like a gate (similar to electronic organs). The poly-mod is modulated exclusively by the assignable envelope. Therefore the poly-mod modulation is completely free at the expense of a less dynamic sound evolution.
\end{itemize}  

\scalebox{0.4}{
  \begin{tikzpicture}[scale=0.8]
   \envroute{0,0}
  \end{tikzpicture}
}

\textbf{Oscillator sync and poly-mod}

As explained in section \ref{osc}, you can sync oscillator A to oscillator B using the \oscsync switch in the sub-panel for oscillator B. This can also be view as a variant of cross-modulation. The sync function can also be combined with frequency modulation as described above. 

Combining oscillator sync with poly-mod pitch modulation by envelope can be used to produce interesting short transients or slow sound evolutions. For example, when oscillator A is a pulse wave, then sync'ing this wave to oscillator B and time modulating the pitch of oscillator A has the effect of pulse width modulation by envelope. It creates a change in timbre which can be short to simulate plucking an brassy attacks or long to produce a sweep of the sound spectrum.

\textbf{Starting points for patch design with poly-mod and sync}

Using poly-mod in combination with oscillator sync, different modulations strengths for different oscillator pitches with different filter and modulation shapes - this offers a very wide field of experimentation. Here are two prototypical starting points:

\begin{itemize}
  \item Set the envelope routing to \textit{standard}, tune both oscillators to the same frequency. In the poly-mod section turn source VCO B to zero and deactivate the filter modulation. Then set VCO A to sync and switch on the poly-mod target frequency oscillator A. Experiment with the poly-mod envelope source amount, oscillator A frequency and filter envelope attack and decay. This produces an evolving sync sound which, for high envelope amounts, can be tuned to a very assertive synth solo sound.
  \item Set the envelope routing to \textit{standard}, tune both oscillators to the same frequency. In the poly-mod section turn the poly-mod envelope source amount to zero, deactivate the filter modulation  and switch on the poly-mod target frequency oscillator A. Switch off VCO A sync. Now experiment with the poly-mod envelope source amount and oscillator B fine tune. Either you enjoy the drastic spectra you will get for most of the settings or you carefully adjust the frequencies to find the sweet spots at which the spectrum clears. This is the typical way to produce very rich and brassy timbres but normally this also needs some filter cut-off to softly cut away the highest frequencies. For stronger poly-mod, oscillator B fine tune may not be sufficient and it may be necessary to also adjust the frequency of oscillator A. For this it is useful to set the granularity of the frequency dials to free/octave (i.e. oscillator A free and oscillator B octave). Notice that other controls also influence the pitch, such as the oscillator B wave selection and pulse width (see description of bleed bug below) so that you may find that you need to compensate for this effect with fine tuning.   
\end{itemize} 

\textbf{Technical excursion: pulse width bleed bug}

The firmware versions 2.0 and 2.1 RC3 contain a software error which incorrectly resets the pulse widths of oscillators A and B. This manifests itself under the following conditions: \polymodosc is turned up and the poly-mod target \polymodfreq is active. The wave shape pulse of oscillator B is not activated. In this case the incorrect reset of the pulse widths leads to a wrong frequency in oscillator A. The effect becomes stronger with higher settings of \polymodosc. In version \version this bug is corrected but in order to ensure backwards compatibility of older patches, the bug is automatically switched on when loading such patches to faithfully reproduce the behaviour of prior firmware versions.  The bug toggle is the \syncbug which can be set to \textit{off} and \textit{on}. The default value (default patch) is \textit{off}.
